\section{Introduction}
\label{sec:introduction}

Parity is a basic test of whether a shallow circuit can combine
information from all of its inputs. For a computational-basis string
$x=(x_1,\ldots,x_n)\in\{0,1\}^n$, its value is
$\parity(x)=x_1\oplus\cdots\oplus x_n$. We study circuits generated
by arbitrary one-qubit unitaries and Toffoli gates with any number of
controls and one target. Gates in each parallel layer have pairwise
disjoint supports. The depth is the number of such layers.

The circuit starts in $\ket{x}\otimes\ket0^{\otimes(N-n)}$ on $N$
total qubits. Its answer is the result of measuring one designated output
qubit in the computational basis; the remaining registers may be
discarded in arbitrary final states, and the input need not be preserved.
For $0<\varepsilon\le1/2$, computing parity with worst-case advantage
$\varepsilon$ means giving the correct answer with probability at least
$1/2+\varepsilon$ for every $x$. The class $\QAC$ uses constant depth
and polynomially many total qubits. This also bounds the gate count:
a disjoint layer on $N$ qubits has at most $N$ nonempty gates.
The gates may depend on $n$, but not on the input string.
We use unitary circuits, with no intermediate measurements,
postselection, or primitive unbounded-fanout gate.

\begin{theorem}[Parity lower bound]
\label{thm:parity}
Fix a nonnegative integer \(d\), a real number \(c\ge1\), and
\(0<\varepsilon\le1/2\). For all sufficiently large \(n\), no circuit
of depth at most \(d\) on \(n\) input qubits and at most \(n^c\) total
qubits, using arbitrary one-qubit unitaries and unbounded-arity Toffoli
gates, computes \(\parity(x)\) with probability at least
\(1/2+\varepsilon\) for every \(x\in\{0,1\}^n\).
Gates in each layer have disjoint supports, ancillas start in \(\ket0\),
and one output qubit is measured; the other qubits may be discarded.
\end{theorem}

Taking \(\varepsilon=1/6\) gives the usual success threshold \(2/3\).
Thus \Cref{thm:parity} resolves positively the conjecture
\(\parity\notin\QAC\), with bounded error, polynomially many ancillas,
and only a measured output required.

\subsection{History and significance}
In the classical class $\mathrm{AC}^0$, circuits use unbounded-fan-in
AND and OR gates and NOT gates, with unrestricted reuse of wires.
The superpolynomial constant-depth lower bounds for parity of Furst,
Saxe, and Sipser~\cite{FurstSaxeSipser1984} and, independently,
Ajtai~\cite{Ajtai1983}, were fundamental explicit circuit lower bounds.
Yao obtained exponential bounds~\cite{Yao1985}, and H\aa stad obtained
almost optimal exponential bounds through the switching
lemma~\cite{Hastad1986}. Quantum circuits introduce a different
constraint: the supports of gates in a parallel layer are disjoint,
so a wire cannot control arbitrarily many simultaneous gates.

Moore introduced the corresponding quantum circuit classes and related
coherent parity to unbounded fanout~\cite{Moore1999}. The latter is the
unitary sending $\ket{b,y_1,\ldots,y_k}$ to
$\ket{b,y_1\oplus b,\ldots,y_k\oplus b}$. The coherent parity
transformation sends the same input to
$\ket{b\oplus\parity(y),y_1,\ldots,y_k}$, where
$y=(y_1,\ldots,y_k)$. Hadamard conjugation interchanges these two
transformations. Moore's original
formulation required ancillary qubits to return to zero; our result
also treats the more permissive measured-output model.
Green, Homer, Moore, and Pollett developed the associated quantum
counting classes and modular-gate equivalences~\cite{GreenHomerMoorePollett2002}.
The constant-depth constructions of H\o yer and \v Spalek demonstrate
how much computational power a primitive fanout operation
provides~\cite{HoyerSpalek2005}.

Early lower bounds revealed the importance of the ancillary register.
Fang, Fenner, Green, Homer, and Zhang obtained a depth lower bound of
order $\log(n/(a+1))$ for exact, clean parity computation with $a$
ancillary qubits~\cite[Corollary 4.6]{FangFennerGreenHomerZhang2006};
Bera developed a different lower-bound method in the ancilla-free
setting~\cite{Bera2011}. Pad\'e, Fenner, Grier, and Thierauf ruled out
exact clean simulation at entangling depth two~\cite{PadeFennerGrierThierauf2020}.
Fenner, Grier, Pad\'e, and Thierauf subsequently removed the final-state
cleanliness requirement for exact single-output parity at that
depth~\cite[Theorem 6.1]{FennerGrierPadeThierauf2025}.
Here the quoted depth-two results count multiqubit layers, allowing
arbitrary one-qubit layers between them; we call this
\emph{entangling depth}. Rosenthal proved lower bounds for approximate
parity unitaries and related cat-like states, and gave constant-depth
approximation constructions with superpolynomial
resources~\cite{Rosenthal2021}. Approximate unitary implementation and
computation of a single measured classical bit have different output
requirements, so the distinction matters in comparing lower bounds.

More recent work developed spectral and Fourier methods for $\QAC$.
Nadimpalli, Parham, Vasconcelos, and Yuen used the Pauli spectrum to prove
ancilla-restricted lower bounds~\cite{NPVY2024}. Anshu, Dong, Ou, and Yao
ruled out a fixed positive parity advantage with
$\widetilde O(n^{1+3^{-d}})$ ancillas at every fixed entangling depth
$d\ge1$~\cite[Corollary 4.4 and Theorem 4.6]{ADOY2025}.
Dong, Ou, and Yao increased this ancilla range to
$\widetilde O(n^{1+2^{-d}})$~\cite[Corollary 3.3]{DOY2025LinearSize}.
Here $\widetilde O$ suppresses logarithmic factors. Both results use
spectral-norm approximation of conjugated observables.

Joshi, Tal, Vasconcelos, and Wright proved exact parity lower bounds at
entangling depth three and bounded parity correlation at entangling
depth two by $\exp(-\Omega(\sqrt n))$, without a size
restriction~\cite[Theorem 1.2 and Corollary 5.9]{JTVW2026}. Kintali's September 2026 preprint reports
correlation at most $\exp(-\Omega(n))$ at entangling depth three
and an exponential multiqubit gate lower bound for exact parity at
entangling depth four~\cite[Theorems 1--2]{Kintali2026ShallowParity}.
These results retain a resource restriction or a fixed upper bound on
entangling depth. Gretta, Gupta, and Joshi established an
equivalence between the general parity lower-bound question and Fourier
concentration~\cite{GGJ2026}, and Xu and Li extended related fanout
reductions to symmetric Boolean functions~\cite{XuLi2026Fanout}. Our theorem permits every fixed depth and
every fixed polynomial bound on the total number of qubits, with
arbitrary discarded garbage and any fixed positive worst-case advantage.

\subsection{The localization argument}

Unbounded gate arity is the first obstacle to a lower bound. A single
Toffoli gate can involve every input, so the number of wires in an
output's backward light cone gives no useful restriction. Our argument
instead bounds transitions from a fixed product condition to states
with many mismatches from that condition. We keep the amplitude lost
when imposing a projection, without normalizing the resulting state.

A \emph{product projection} fixes a chosen one-qubit vector on each qubit
of a subset and acts as the identity elsewhere. Fix reference unit
vectors $v_i$ on a subset $S\subseteq\{1,\ldots,N\}$. Their
\emph{mismatch count} is
\[
 M=\sum_{i\in S}(I-\proj{v_i})_i;
\]
its eigenvalue counts the factors orthogonal to their reference vectors.
Write $[M=0]$ and $[M\ge r]$ for its zero and high-count spectral
projections. Given a shallow circuit $U$ and a product projection $A$,
our structural theorem bounds
\[
 \norm{[M\ge N^t]UAU^\dagger[M=0]}\le e^{-N^s}
 \qquad(0\le s<t)
\]
for every fixed depth and all sufficiently large $N$. The norm is the
operator norm, and the estimate is uniform in $U,A,M$. In particular,
the uncounted register can carry an arbitrary, internally entangled state.
\Cref{thm:localization} gives the precise statement in a reflection-gate
model containing the original circuits.

There are two inductions. The outer induction removes one circuit layer.
Every allowed Toffoli is a reflection $I-2P$ about a product projection.
If $A$ fixes only $k$ qubits, at most $k$ gates in a disjoint layer meet
its support. Expanding those reflections on the two sides of $A$
expresses the conjugate as a sum of products of three projections,
each conjugated by a circuit of one smaller depth. Intermediate mismatch
cutoffs bound these ordered products. This proves localization for
projections supported on few qubits (\Cref{sec:depth}).

The reflection expansion has a cost exponential in the support size,
so it does not directly control arbitrary product projections. To handle
such a projection, write it as $[D=0]$ for a second mismatch count $D$.
A polynomial $p(D)$ of low degree expands
into projections on few qubits. The construction in \Cref{sec:polynomial}
controls both its approximation error on a low-count interval and the
sum of its coefficients weighted by the cost of this expansion.
It need not be a good approximation at large counts. The second
induction controls precisely that remaining error.

Here the order of the operators is decisive. Put $P=[M=0]$,
$Q=U[D=0]U^\dagger$, and $H=[M\ge N^t]$. Squaring the transition
norm gives $\norm{HQP}^2=\norm{HQPQH}$. Replacing only the first
$Q$ by $p(UDU^\dagger)$ leaves an error acting on $PQ$. After
conjugation by $U^\dagger$, the high-count part of that error is another
localization problem, with the two counts exchanged and a larger
mismatch threshold. We prove the universal statement at that larger
threshold first. A finite chain of thresholds ends above the maximum
possible count $N$, where localization is automatic; working backward
proves the original statement (\Cref{sec:bootstrap}). The same-depth
inverse-circuit invocation is therefore part of a finite induction.

Finally, the output-zero probability is the diagonal of the conjugated
output projection. Its parity coefficient is the average of matrix
entries between complementary Hadamard product states. Those states
differ on all $n$ input qubits. The localization theorem makes every
such entry small when $N$ is polynomial in $n$, whereas a fixed
worst-case advantage forces their average to stay positive
(\Cref{sec:parity}). This final step places no restriction on the
unmeasured registers.

The polynomial ingredients have a substantial history. The combination
of Chebyshev approximation with exact roots at exceptional integer
points follows an established approximation-and-interpolation pattern
for symmetric predicates~\cite[Section 3.1]{Sherstov2009}.
Polynomial localization and concentration for quantum states appear in
work of Kuwahara, Arad, Amico, and
Vedral~\cite{KuwaharaAradAmicoVedral2017} and Anshu and
Metger~\cite{AnshuMetger2023}. The latter also track the total norm of
local terms. For $\QAC$, the spectral-norm approximation approach of
Anshu, Dong, Ou, and Yao~\cite{ADOY2025} is a particularly relevant
predecessor. We prove the required polynomial and coefficient estimates
here. The finite inverse-circuit induction supplies the additional scale
improvement: localization holds for every fixed exponent pair
$0\le s<t$, uniformly in both counts and the entire uncounted register.

The companion article \emph{Regular trajectories, pruning and quantum
parity}~\cite[Sections 4--7]{OpenAI2026Trajectories} gives an independent proof through
state-dependent propagation, projection insertion, and pruning.
Its structural estimates concern regular trajectories; the present
operator-norm theorem is proved without a trajectory hypothesis.
Both articles prove their own polynomial estimates and their own parity
conversion, so neither parity proof invokes the other.

\subsection{Consequences}

Localization and parity also constrain symmetric computation and state
preparation. We combine them with reductions of Xu and Li~\cite{XuLi2026Fanout}
and state constructions of Gretta, Gupta, and Joshi~\cite{GGJ2026}.
In particular, strict majority,
$\operatorname{MAJ}_n(x)=\mathbf 1\{|x|>n/2\}$, where $|x|$ is the
Hamming weight, cannot be computed in this circuit model with worst-case advantage $1/(\log n)^\gamma$
for any fixed $\gamma>0$ (\Cref{cor:symmetric-functions}). This separates
the model from the corresponding one with unbounded threshold gates.

The state-preparation consequences rule out inverse-polynomial
measurement probabilities for both a string and its bitwise complement,
and sufficiently accurate preparation of uniform superpositions of
strings of a fixed Hamming weight $k$ with $n^\delta\le k\le n/2$
for fixed $\delta>0$ (\Cref{cor:state-preparation}). For a single measured output, define
its expected sign as the probability of zero minus the probability of one.
Localization implies that its Fourier coefficients have total squared
weight that decays faster than every inverse power of $n$ on input
subsets of size at least $n^\delta$, for every fixed
$0<\delta\le1$ (\Cref{cor:quantum-lmn}). Finally, majority cannot be computed
with average success at least $1-\tfrac12 n^{-\delta}$ on uniformly
random inputs, for any fixed $\delta>0$ (\Cref{cor:majority-average}).
The precise quantifiers and proofs
appear after the parity proof in \Cref{sec:consequences}.

\paragraph{Organization.}
\Cref{sec:localization} states the localization theorem and fixes the
circuit model. \Cref{sec:polynomial} proves the polynomial approximation.
\Cref{sec:depth} reduces small-support projections to the preceding
depth, and \Cref{sec:bootstrap} completes the finite scale induction.
\Cref{sec:parity} derives \Cref{thm:parity}, and
\Cref{sec:consequences} proves the consequences just described.
